% Machine-Enforced Research Governance for Quantitative Finance
% Target: arXiv q-fin.CP preprint -> Journal of Financial Data Science
% Build: make (latexmk). Claims discipline: every numeric claim in this
% document must exist in paper/claims.json and be re-derived by
% tests/unit/test_paper_claims.py from committed research artifacts.
\documentclass[11pt]{article}

\usepackage[margin=1.1in]{geometry}
\usepackage{amsmath,amssymb,amsthm}
\usepackage{booktabs}
\usepackage{graphicx}
\usepackage[hidelinks]{hyperref}
\usepackage[round]{natbib}
\usepackage{xcolor}

\newtheorem{definition}{Definition}
\newcommand{\todoA}[1]{\textcolor{red}{[PENDING #1]}}
\newcommand{\CONFIRMED}{\textsc{confirmed}}
\newcommand{\HYPOTHESIS}{\textsc{hypothesis}}

\title{Machine-Enforced Research Governance for Quantitative Finance:\\
A Production System Where Pre-Registration, Multiple-Testing\\
Control, and Out-of-Sample Validation Are Not Optional}

\author{Tim Wu\thanks{Single human author; this work was carried out with
substantial AI assistance, itself an object of study --- see the AI
Involvement Disclosure (Section~\ref{sec:ai}).}}

\date{Draft v0.1 --- \today}

\begin{document}
\maketitle

\begin{abstract}
The empirical finance literature has diagnosed its own replication crisis
and prescribed remedies: multiple-testing correction
\citep{harvey2016and}, pre-registration \citep{harvey2017presidential},
and full disclosure of all trials \citep{bailey2017probability,
fabozzi2018being}. A decade later, every one of these remedies remains
voluntary. We describe, to our knowledge, the first \emph{publicly
documented and open-source} production system in which these norms are
enforced by software rather than by goodwill: every statistical
regularity must pass a false-discovery-rate gate \emph{and} a
walk-forward out-of-sample gate --- treated by the closest prior work as
alternatives \citep{harvey2016and} --- before it may be labeled
\CONFIRMED{}; every component change follows a frozen, git-tagged
pre-registered protocol with numeric decision gates fixed before any test
code exists; every candidate must pass a planted-truth positive-control
power check before touching real data; failures are recorded in permanent
public ledgers with the same ceremony as successes; a user-facing
large-language-model agent is bound to disclose machine-readable evidence
grades and data provenance; and an autonomous AI research agent operates
under the same governance with all promotion triggers reserved to a
human. We report a pre-registered 50-year, five-cohort validation study
of the underlying system, a public scoreboard of ten pre-registered
method tests (two promoted, eight fast-failed), three in-vivo case
studies in which the governance caught real defects --- including a
calibration bug in our own Bayesian instrument --- and a counterfactual
ablation quantifying what weaker governance regimes would have wrongly
adopted. All code, protocols, ledgers, and results are public and
bit-for-bit reproducible.
\end{abstract}

% =====================================================================
\section{Introduction}\label{sec:intro}

Quantitative finance has a peculiar relationship with its own
epistemology. The field's leading journals have published the diagnosis
of a replication crisis in unusual detail: most claimed factor
discoveries are likely false under any defensible multiple-testing
threshold \citep{harvey2016and}; backtest overfitting is nearly
guaranteed when many configurations are tried on the same data
\citep{bailey2017probability}; and the incentive structure rewards
publishing fragile positives \citep{harvey2017presidential}. The
prescriptions are equally well documented: raise discovery thresholds
and correct for multiplicity \citep{harvey2016and}, pre-register
hypotheses and analysis plans before touching data
\citep{harvey2017presidential, chambers2013registered,
bloomfield2018gathering}, and disclose every trial behind a reported
result \citep{fabozzi2018being}.

What the literature does not contain --- so far as an extensive
adversarially-verified search could establish (Section~\ref{sec:related})
--- is an account of these prescriptions being \emph{enforced}: encoded
as hard gates in a running system, such that an un-corrected,
un-validated, or un-registered claim is not merely frowned upon but
mechanically impossible to promote. This paper describes such a system
in production, and treats three questions:

\begin{enumerate}
\item \textbf{Enforceability.} Can the voluntary norms of the
  replication-crisis literature be translated into machine-checkable
  gates without paralyzing research? We give a constructive answer: a
  protocol lifecycle (draft $\to$ power-gated dry run $\to$ freeze $\to$
  mechanical verdict $\to$ ledger) that ran ten complete cycles.
\item \textbf{What enforcement buys.} Does the machinery actually block
  anything a competent human would not? Our public scoreboard --- two
  promotions, eight fast failures, several of them literature-celebrated
  methods --- and a counterfactual ablation (Section~\ref{sec:ablation})
  quantify the answer.
\item \textbf{AI under governance.} As LLM agents enter the research
  loop, who checks their claims? We bind both a user-facing agent
  (machine-readable evidence labels, mandatory provenance disclosure)
  and an autonomous research agent (power-gated screening, human-only
  promotion triggers) to the same rules, and evaluate compliance
  empirically (Section~\ref{sec:agenteval}).
\end{enumerate}

Our claims are deliberately bounded. The mathematical components of the
system are mature, cited methods --- that is a design choice, not an
accident (Section~\ref{sec:system}). The contribution is the governance
layer and the empirical demonstration that it works end to end. And
because industry systems are rarely published, we claim only the first
\emph{publicly documented} implementation, not the first implementation.

% =====================================================================
\section{Related Work}\label{sec:related}

\paragraph{The replication crisis and its prescriptions.}
\citet{harvey2016and} re-examined 316 published factors and concluded
that a $t$-statistic threshold of roughly 3.0 or more is required to
control false discoveries; crucially for our purposes, that paper
discusses out-of-sample validation and multiple-testing frameworks as
\emph{alternative} remedies and adopts the latter, noting that
out-of-sample tests ``cannot be used in real time.''
\citet{bailey2017probability} formalized the probability of backtest
overfitting via combinatorially symmetric cross-validation;
\citet{fabozzi2018being} proposed a voluntary disclosure template for
multiple testing; \citet{harvey2017presidential} recommended
pre-registration in finance, citing the Open Science Framework, and the
2017 \emph{Journal of Accounting Research} conference implemented
registered reports \citep{bloomfield2018gathering}. Our system differs
from this entire line in mechanism, not in motivation: what these works
recommend to humans, our system refuses to let humans skip.

\paragraph{The components' own prior art.} Regime-switching models in
finance descend from \citet{hamilton1989new}; see
\citet{ang2012regime, nystrup2018dynamic}. Transfer entropy
\citep{schreiber2000measuring} applied to finance dates to
\citet{marschinski2002analysing}, is Granger-equivalent under
Gaussianity \citep{barnett2009granger}, has standard financial-market
applications \citep{dimpfl2013using}, and is available with FDR-aware
inference in open tooling \citep{wollstadt2019idtxl}; Granger-causal
network scans of financial series are standard since
\citet{billio2012econometric}. Local projections are due to
\citet{jorda2005estimation}, with Bayesian variants in
\citet{ferreira2023bayesian}. Volatility-filtered extreme value theory
for conditional tail risk is \citet{mcneil2000estimation}.
Walk-forward validation was codified by \citet{pardo1992design}. We
claim novelty for none of these.

\paragraph{What we could not find.} An adversarially verified literature
search (109 research agents, three-vote verification per claim; search
artifacts in the repository) located no prior work in which (i)
FDR control \emph{and} walk-forward out-of-sample validation are jointly
necessary conditions enforced by software, (ii) the resulting binary
epistemic label is machine-readable and binds a user-facing LLM agent's
disclosures, or (iii) an autonomous research agent is subject to
pre-registration and power gates with human-only promotion. Absence of
evidence is not evidence of absence --- proprietary systems may do some
of this unpublished --- hence our bounded claim.

% =====================================================================
\section{The Underlying System}\label{sec:system}

PRISM (Probabilistic Regime \& Shock-Impact Model) is a nightly batch
that fits, on 10--50 years of public macro-financial data
(Section~\ref{sec:repro}), four component families, each mature:
a Gaussian hidden Markov regime model \citep{hamilton1989new};
a causality scan combining transfer entropy
\citep{schreiber2000measuring, marschinski2002analysing} and Granger
tests under Benjamini--Hochberg FDR control
\citep{benjamini1995controlling} with Fisher combination
\citep{fisher1925statistical} across horizons; Bayesian local
projections \citep{jorda2005estimation, ferreira2023bayesian} with a
conjugate normal-inverse-gamma prior; and volatility-filtered
generalized-Pareto tail estimation \citep{mcneil2000estimation},
backtested with standard coverage tests \citep{kupiec1995techniques,
christoffersen1998evaluating}. Outputs are served to a user-facing LLM
agent that answers questions about markets, regimes, and risk.

The deliberate maturity of the components is load-bearing for this
paper's argument: whatever value the system demonstrates comes from
\emph{how claims are admitted}, not from novel estimators.

% =====================================================================
\section{The Governance Layer}\label{sec:governance}

\subsection{Epistemic labels and the dual gate}\label{sec:dualgate}

Every statistical regularity the system may assert carries exactly one
of two labels.

\begin{definition}[\CONFIRMED{}]
A regularity is \CONFIRMED{} if and only if (a) its causality-scan
$q$-value passes the pre-registered FDR threshold, \emph{and} (b) it
exhibits a positive out-of-sample sign-accuracy edge over an AR(1)
baseline in a pre-specified fraction of expanding-window walk-forward
folds. Otherwise it is \HYPOTHESIS{}, permanently and by default.
\end{definition}

The conjunction is the point. The closest prior work treats (a) and (b)
as substitutes \citep{harvey2016and}; our position, learned from
operating the system, is that they fail differently --- FDR control does
not stop a regularity that is real in-sample but unstable, and OOS
validation alone is silent about the multiplicity of the search that
produced the candidate --- so only the conjunction is safe to expose to
downstream consumers. The label is data, not prose: it is stored with
the regularity, served through the API, and the user-facing agent's
system prompt hard-requires quoting it (evaluated in
Section~\ref{sec:agenteval}).

\subsection{Data-provenance labels}

Every quantity the agent may quote carries one of four provenance
classes: \textsc{live} (licensed market data), \textsc{derived}
(computed signals, hypothesis-grade), \textsc{model} (PRISM outputs,
carrying the Section~\ref{sec:dualgate} label), or \textsc{simulated}
(demonstration data). Silent mixing is prohibited by prompt contract
and audited empirically (Section~\ref{sec:agenteval}).

\subsection{The protocol lifecycle}\label{sec:lifecycle}

Component changes and candidate methods follow a state machine:

\begin{enumerate}
\item \textbf{Draft.} A protocol file states the null hypothesis
  (always ``the candidate does \emph{not} help''), the exact incumbent,
  numeric decision gates, cohort definitions, seeds, known limitations,
  and a fast-failure clause. Drafts carry a header declaring they confer
  no confirmatory status.
\item \textbf{Dry run with power gates.} Before freezing, every
  instrument must detect a planted, literature-calibrated effect in
  synthetic data with power $\geq 0.80$ at controlled false-positive
  rate, exercised \emph{through the production entry point}
  (Section~\ref{sec:power}). Designs whose constraints are jointly
  unsatisfiable die here, before consuming their one shot at real data.
\item \textbf{Freeze.} The protocol is committed and git-tagged before
  any real-data test code exists. Moving gates after the tag voids the
  confirmatory claim by construction.
\item \textbf{Mechanical verdict.} Gates are evaluated as written. A
  significant result in the wrong direction is evidence of harm, not
  nuance.
\item \textbf{Ledger.} Outcomes are committed regardless of sign:
  promotions to the production system; failures to a permanent
  scoreboard; pre-freeze screening kills to a rejected-candidates
  ledger; representation-bound deferrals (methods right for a scale the
  system does not yet have) to a conditional-candidates ledger with
  explicit re-activation triggers.
\end{enumerate}

Re-tests are one-shot: a failure caused by recorded protocol defects may
justify a v2 with explicit human sign-off, fixing only the recorded
defects; a second failure of any kind closes the candidate terminally.

\subsection{An autonomous research agent under the same rules}
\label{sec:scout}

A scheduled LLM research agent executes only the first two lifecycle
stages: it reads the ledgers (so dead candidates cannot be re-shopped),
surveys the literature against the system's current constraint
fingerprint, runs the planted-truth power screen on at most one
candidate per cycle, and opens a pull request containing either a draft
protocol or a one-line rejection-ledger entry. It cannot freeze, tag,
merge, or touch production paths; every promotion trigger is a human
action. Its first production cycle rejected a literature-celebrated
candidate (208 citations) whose power screen topped out at 28.7\%
against the 80\% gate --- and, in the same run, surfaced the calibration
defect described in Section~\ref{sec:cases}.

% =====================================================================
\section{Validation Methodology}\label{sec:method}

This section answers, in first-principles terms, how the tests that
back every claim in Section~\ref{sec:results} were built --- because a
result is only as credible as the instrument that produced it.

\subsection{Planted-truth calibration of every instrument}

No statistical tool is used on real data before demonstrating, on
synthetic data with known ground truth, that it (i) finds a planted
effect of the documented size, (ii) stays silent on pure noise at its
nominal false-positive rate, and (iii) is exercised through the exact
code path production uses --- not an internal function with
hyperparameters bypassed. Requirement (iii) exists because we were
burned without it: an early calibration called a core statistic with a
fixed lag while the production path selected the lag by BIC, silently
destroying the mechanism the test certified.

\subsection{Positive-control power gates}\label{sec:power}

\begin{definition}[Power gate]
Let $T$ be a candidate test with decision gate $G$, and let $\theta^*$
be the smallest effect the literature documents for the candidate
mechanism. Simulate $S$ datasets from a null-faithful generator with
$\theta^*$ planted, using the \emph{real} exogenous regressor path
where one exists (its autocorrelation, not the sample size, typically
binds power). The candidate may proceed to real data only if
$\Pr[G \text{ passes}] \geq 0.80$ under the plant and
$\Pr[G \text{ passes}] \leq \alpha$ under the null.
\end{definition}

Two operational lessons are codified. First, constraints must be checked
\emph{jointly}: one campaign candidate fixed a frequency-resolution
defect and then died because the fix left no statistical power ---
resolution and power were jointly unsatisfiable, discoverable in one
simulation. Second, the gate prunes honestly in both directions: in the
study of Section~\ref{sec:enso}, the same simulation that passed a
producer-price endpoint at power 1.00 excluded a consumer-price endpoint
at power 0.34, recording it as untestable rather than testing it anyway.

\subsection{Point-in-time discipline and mutation checks}

At date $t$, instruments may read only data available at $t$; smoothed
full-sample outputs (e.g., HMM posterior probabilities) are wrapped in
expanding-window refits taking the filtered last observation, and the
gap between smoothed and point-in-time scores is reported as the
look-ahead bias it is. Test suites are mutation-checked: for each guard
we verify that re-introducing the defect it protects against makes the
test fail, which kills vacuous assertions before they certify anything.

% =====================================================================
\section{Empirical Results}\label{sec:results}

\subsection{A pre-registered 50-year validation study}\label{sec:study}

The system itself was validated under its own rules before any method
tests ran: protocol frozen at tag \texttt{study-preregistered}, then
back-to-back evaluation on five cohorts (C50: 1976--2026 through C10:
2016--2026) against eight hypothesis families. \todoA{claims.json:
verdict matrix import + headline numbers (AUROC 0.84--0.86 vs NBER;
VaR calibration; Hansen SPA nulls) with exact artifact paths}
The study's headline honesty results: the regime layer is a valid
historical lens but lags real-time detection by more than ten business
days and is therefore served as context, never as an early warning; and
no economic predictability survived superior-predictive-ability testing
\citep{hansen2005test}, which is disclosed to users as a hard rule
against crash-timing claims.

\subsection{The ten-test scoreboard}\label{sec:scoreboard}

\begin{table}[h]
\centering\small
\begin{tabular}{llll}
\toprule
\# & Candidate & Verdict & Gate that decided \\
\midrule
1 & Regime-conditional EVT (own design) & fast fail & regime lag compounds at transitions \\
2 & Fisher causality combination & \textbf{promoted} & power 10--55\%$\to$85--100\% \\
3 & Vol-filtered EVT VaR \citep{mcneil2000estimation} & \textbf{promoted} & calibration green, both assets \\
4 & QIS covariance shrinkage & fast fail & no help at $p{=}5$ (conditional ledger) \\
5 & RFSV rough volatility & fast fail & daily proxy loses CC parity \\
6 & Signature (L\'evy-area) features & fast fail & 0/3 cohorts; AUROC $-0.10$ \\
7 & BOCPD/CUSUM fast alarm & fast fail & + 2 protocol defects recorded \\
8 & BOCPD v2 (one-shot re-test) & terminal & sign-permutation $p=0.172$ \\
9 & Breitung--Candelon frequency causality & fast fail & protocol design (lag collapse) \\
10 & BC v2 (one-shot re-test) & terminal & positive control invisible on real data \\
\bottomrule
\end{tabular}
\caption{Public scoreboard of pre-registered method tests. Two
promotions, eight fast failures; the failures include
literature-celebrated methods, and the ratio is the methodology's
credibility --- a gate that never rejects is not a gate.}
\label{tab:scoreboard}
\end{table}

Note that own designs and external candidates face identical gates:
entry 1 is our own idea, refuted with the same ceremony as the rest.

\subsection{Counterfactual ablation: what weaker governance adopts}
\label{sec:ablation}
\todoA{A1 experiment: pre-registered protocol; replay stored campaign
results under (i) na\"ive in-sample $p<0.05$, (ii) FDR-only, (iii)
OOS-only, (iv) dual gate; adoption matrix table}

\subsection{Agent disclosure compliance}\label{sec:agenteval}
\todoA{A2 experiment: pre-registered protocol; $N\geq50$ live trials
$\times$ prompt battery; label-compliance rates with Wilson intervals}

\subsection{Three in-vivo case studies}\label{sec:cases}

\paragraph{The governance catching our own instrument.} The autonomous
agent's first cycle flagged that the local-projection prior looked
mis-scaled. A planted-truth audit through the production entry point
confirmed it: the absolute-scale inverse-gamma prior contributed 88.7\%
of the posterior scale at horizon 1 on daily-return-magnitude data,
inflating 95\% credible bands 2.96$\times$ (h=1), 1.61$\times$ (h=5),
1.19$\times$ (h=20), with empirical coverage of 100\% against a nominal
95\%. The fix (scale-matched weak prior) restored bands to within
1--2\% of the frequentist oracle; an independent review of the fix then
found that the old prior had silently served as a degenerate-data
floor, and the repair of that regression was itself mutation-checked.
Post-deployment, production band-widths narrowed by the predicted
factors (1.98$\times$ at h=1 for a volatility-index/equity pair) ---
theory, simulation, and production agreeing is precisely what the
methodology is for. \todoA{claims.json: wire each number to
findings\_addendum + PR artifacts}

\paragraph{An honest rejection.} Smooth Local Projections --- an exact
deficiency match with high literature standing --- was rejected before
any real-data spend when its positive-control power reached only 28.7\%
against the 80\% gate under the system's dimensionality, and the
deficiency it targeted was later traced to the instrument defect above:
the ``missing method'' was a calibration bug.

\paragraph{Joint-constraint discovery in gate design.} In the
climate-transmission study (Section~\ref{sec:enso}), the production
walk-forward rule proved degenerate at monthly frequency and
\emph{no} out-of-sample rule achieved 80\% power at the 12-month
horizon (roughly 50 independent draws in 50 years); a calibrated
pooled-edge rule at the 6-month horizon reached power 0.817 at
false-positive rate 0.05. The gate was redesigned before freezing ---
visible, recorded, and legitimate precisely because the protocol was
still a draft.

\subsection{A worked pre-registration: climate transmission}
\label{sec:enso}
\todoA{summarize DRAFT\_protocol\_enso\_farmppi.md + power tables as the
end-to-end worked example; note freeze/test pending}

% =====================================================================
\section{Limitations}\label{sec:limits}

(i) This is an $n{=}1$ system study with a single maintainer; the
governance has not been stress-tested against adversarial researchers or
team-scale incentives. (ii) ``First publicly documented'' is a ceiling:
proprietary systems may implement similar controls unpublished. (iii)
The sealed 50-year verdicts were produced under the pre-2026-08-12
instrument; re-runs under the corrected prior are new evidence, not
reproductions, and the repository records this discontinuity explicitly.
(iv) The methodology's domain-generality is itself labeled
\HYPOTHESIS{}: it is distilled from one finance system at daily/monthly
frequency, and its first out-of-domain application is the designated
test. (v) Governance costs are real (ten cycles produced two adoptions);
we argue the eight documented failures are the product, not the waste,
but that argument is philosophical, and the ablation
(Section~\ref{sec:ablation}) is its only quantitative support.

\section{Future Work}\label{sec:future}
Second-domain transfer of the protocol lifecycle; community replication
of the scoreboard; integration with journal registered-reports workflows
\citep{bloomfield2018gathering}; scale-up of the conditional-candidates
ledger as the system's representation grows; and richer compliance
auditing of LLM agents as they take on more of the research loop.

\section{AI Involvement Disclosure}\label{sec:ai}
\todoA{A4: precise account --- system and paper built with an AI
assistant (Claude); the autonomous scout is an AI agent; all promotion
triggers, merges, freezes, deployments, and this paper's claims were
human-gated; drafted with A2 results}

\section{Reproducibility}\label{sec:repro}
All protocols, ledgers, runners, and results are in the public
repository. \texttt{make study} reproduces the 50-year study bit-for-bit
(fixed seeds, cohort checkpoints). Data are public series (FRED,
NOAA/CPC) referenced by identifier; retrieval instructions are committed
alongside each protocol. \todoA{provenance table: series ids, access
dates, commit hashes}

\bibliographystyle{plainnat}
\bibliography{bib/references}

\end{document}
